\documentclass{article}
\usepackage[utf8]{inputenc}
\usepackage{amsmath}

\title{Sample Project}
\author{Your Name}
\date{\today}

\begin{document}
\maketitle

\begin{abstract}
This is a sample project so you can see Local Overleaf working right away.
Edit this file, hit Compile, and watch the PDF preview on the right update.
\end{abstract}

\section{Introduction}
A few things to try:
\begin{itemize}
  \item \textbf{Compile} (top right) to build this PDF.
  \item Click anywhere in the PDF preview to jump to the matching line in this
    source file (SyncTeX), and vice versa with \emph{Show in PDF}.
  \item Open \textbf{History} to see automatic snapshots after every
    successful compile, plus manual named saves.
  \item Toggle \textbf{Dark mode} from the toolbar for long writing sessions.
  \item Use the \textbf{Files} panel to add \texttt{.bib} files, images, and
    folders, or \textbf{Import from Overleaf} on the dashboard to bring in an
    existing Overleaf project via its git URL.
\end{itemize}

\section{Math, since this is a LaTeX editor}
Euler's identity:
\[
  e^{i\pi} + 1 = 0
\]

\section{Conclusion}
Delete this project whenever you're ready to start your own paper.

\end{document}
